\chapter{アルゴリズムの基礎}
\label{ch:seminar-02}

%% ============================================================
%% 線形計画としての定式化
%% ============================================================
\section{線形計画としての定式化}
\label{sec:sem-lp-standard}

離散 Kantorovich 問題
\[
  \MKD_{\mathbf{C}}(\mathbf{a}, \mathbf{b})
  \defeq \min_{\mathbf{P} \in \CouplingsD(\mathbf{a}, \mathbf{b})}
  \inner{\mathbf{C}}{\mathbf{P}}
\]
は線形計画である．
$\mathbf{P}$ と $\mathbf{C}$ をそれぞれベクトル化した
$\mathbf{p}, \mathbf{c} \in \R^{nm}$，
周辺条件をまとめた制約行列 $\mathbf{A} \in \R^{(n+m) \times nm}$ を用いると，
標準形
\[
  \MKD_{\mathbf{C}}(\mathbf{a}, \mathbf{b})
  = \min_{\substack{\mathbf{p} \in \R_{\geq 0}^{nm} \\
    \mathbf{A}\mathbf{p} = [\mathbf{a}^\top, \mathbf{b}^\top]^\top}}
  \mathbf{c}^\top \mathbf{p}
\]
と表される．$nm$ 個の変数と $n + m$ 本の等式制約を持つ有限次元の線形計画であり，
\textbf{ネットワーク単体法}（network simplex）をはじめとする
LP アルゴリズムで厳密に解ける．

%% ============================================================
%% 主双対最適性条件
%% ============================================================
\section{主双対最適性条件}
\label{sec:sem-primal-dual}

線形計画の双対理論から，最適解を特徴付ける主双対条件が得られる．

\begin{claim}{主双対最適性}{sem-primal-dual-optimality}
  $\mathbf{P} \in \CouplingsD(\mathbf{a}, \mathbf{b})$ と
  $(\mathbf{f}, \mathbf{g}) \in \R^n \times \R^m$ が以下を満たすとする：
  \begin{enumerate}
    \item[(i)] \textbf{相補性}：$P_{i,j} > 0 \Longrightarrow f_i + g_j = C_{i,j}$，
    \item[(ii)] \textbf{双対実行可能性}：$f_i + g_j \leq C_{i,j}$（$\forall i, j$）．
  \end{enumerate}
  このとき $\mathbf{P}$ は離散 Kantorovich 問題の最適解であり，
  \[
    \MKD_{\mathbf{C}}(\mathbf{a}, \mathbf{b})
    = \inner{\mathbf{C}}{\mathbf{P}}
    = \inner{\mathbf{f}}{\mathbf{a}} + \inner{\mathbf{g}}{\mathbf{b}}.
  \]
\end{claim}

\begin{proof}
  任意の $\mathbf{Q} \in \CouplingsD(\mathbf{a}, \mathbf{b})$ に対し，
  双対実行可能性 (ii) と $Q_{i,j} \geq 0$ から
  \[
    \inner{\mathbf{C}}{\mathbf{Q}}
    = \sum_{i,j} C_{i,j} Q_{i,j}
    \geq \sum_{i,j} (f_i + g_j) Q_{i,j}
    = \inner{\mathbf{f}}{\mathbf{a}} + \inner{\mathbf{g}}{\mathbf{b}}.
  \]
  $\mathbf{P}$ については相補性 (i) から $P_{i,j} > 0$ の項で等号が成り立ち，
  $P_{i,j} = 0$ の項は和に寄与しないから
  \[
    \inner{\mathbf{C}}{\mathbf{P}}
    = \inner{\mathbf{f}}{\mathbf{a}} + \inner{\mathbf{g}}{\mathbf{b}}.
  \]
  以上から
  $\inner{\mathbf{C}}{\mathbf{P}} \leq \inner{\mathbf{C}}{\mathbf{Q}}$
  が任意の $\mathbf{Q}$ で成立し，$\mathbf{P}$ は最適．
\end{proof}

ネットワーク単体法はこの条件に基づき，
実行可能解を改善しながらコストを単調に減少させ，
多項式時間で最適解に到達する（Orlin, 1997）．

%% ============================================================
%% エントロピー正則化への動機付け
%% ============================================================
\section{エントロピー正則化への動機付け}
\label{sec:sem-motivation-entropic}

ネットワーク単体法は厳密解を与えるが，
大規模応用の観点では根本的な制約をもつ：

\begin{itemize}
  \item \textbf{計算量}：$O(nm)$ 個の変数に対する反復・データ構造管理は，
    $n, m$ が $10^5$ オーダーになると現実的でない．
    GPU 並列化にも本質的に向かない（離散探索的）．
  \item \textbf{微分不可能性}：最適値 $\MKD_{\mathbf{C}}(\mathbf{a}, \mathbf{b})$ は
    $\mathbf{a}, \mathbf{b}$ について区分線形にしかならず，
    自動微分による勾配最適化に組み込みにくい．
\end{itemize}

これらの困難は，目的関数にエントロピー項を加えて
正則化することで大きく緩和される．
正則化された問題は \textbf{Sinkhorn アルゴリズム}で効率的に解け，
計算は GPU 並列に適し，最適値は周辺分布に関して微分可能となる．
